% Generated draft: requires mathematical and source review.
\documentclass[11pt]{article}
\usepackage[margin=25mm]{geometry}
\usepackage{fontspec}
\setmainfont{Latin Modern Roman}
\usepackage{amsmath,amssymb,mathtools,unicode-math}
\setmathfont{Latin Modern Math}
\usepackage{enumitem,needspace,xurl,hyperref}
\setlength{\parindent}{0pt}
\setlength{\parskip}{5pt}
\setlength{\emergencystretch}{3em}
\newcommand{\N}{\mathbb{N}}
\newcommand{\Z}{\mathbb{Z}}
\newcommand{\Q}{\mathbb{Q}}
\newcommand{\R}{\mathbb{R}}
\newcommand{\C}{\mathbb{C}}
\newcommand{\F}{\mathbb{F}}
\newcommand{\A}{\mathbb{A}}
\newcommand{\PP}{\mathbb{P}}
\newcommand{\E}{\mathbb{E}}
\newcommand{\eps}{\varepsilon}
\newcommand{\dd}{\,\mathrm{d}}
\begin{document}
\section*{UnsolvedMath ID: 800477}
\textbf{Source problem code:} catalogue entry
\subsection{Definitions and mathematical statement}
\textbf{Scope and terminology.} The objects under discussion are tiling billiards on periodic tilings of the plane. Under the usual doubly periodic interpretation, a planar tiling \(\mathcal T\) is periodic if there are linearly independent vectors \(v_1,v_2\in\R^2\) such that translation by every element of \(\Z v_1+\Z v_2\) preserves \(\mathcal T\). The supplied excerpt does not specify the admissible tiles, the billiard propagation rule, the treatment of trajectories meeting vertices, or the precise intended periodicity convention. These must be fixed before a complete dynamical formulation is possible.

\textbf{Question attributed to Olga Paris-Romaskevich.} In the class of tiling billiards on periodic planar tilings, is there a finite list of generic archetypal behaviors? In particular, does every such billiard exhibit one of the following two generic behaviors?
\begin{enumerate}
\item Typical trajectories are dense in open sets of the plane.
\item The behavior is integrable, in the sense used in the source: typical trajectories are either periodic or linearly escaping.
\end{enumerate}
The second question is a proposed dichotomy within the broader finite-classification question. Here ``generic,'' ``typical,'' ``linearly escaping,'' and ``dense in open sets'' retain their source meanings, which are not defined in the supplied excerpt. In particular, no measure, topology, trajectory space, or exceptional-set convention is supplied. Consequently, the displayed alternatives cannot yet be interpreted as a fully specified quantified assertion about trajectories, and the allowed meaning of an archetypal behavior remains unresolved.

\textbf{Supplied context and qualifications.} The source reports the first alternative for its cases 1 and 5, and integrable behavior for its cases 2 and 3; integrability in case 4 is explicitly conjectural. Those cases are not described in the supplied record.

The source also reports that locally foldable periodic tiling billiards are integrable subject to technical regularity conditions. Neither local foldability nor those conditions are specified here, so this contextual claim must not be read as an unconditional theorem in this draft. The source attributes the argument to fine Morse theory developed by Ivan Dynnikov, particularly Theorem 1 of [Dynnikov\_main], together with [PaRoH]. It describes an ingredient originating in [BDFI18] and developed further in [TAF, PaRoH]: a locally foldable tiling admits a global folding map \(\mathcal F\) on the tiled plane whose restriction to neighboring tiles folds one onto the other along their common edge by axial reflection. The precise hypotheses and construction are not supplied.

The excerpt proposes irrational brick tiling billiards and the parallelogram tiling billiards discussed in its Question 2 as possible initial tests among simple tilings that are not locally foldable. Their definitions are absent. Its reference to testing ``Question 1'' is retained as an unresolved internal cross-reference; this registry record is attributed to Problem 3.
\subsection{Short English statement}
Olga Paris-Romaskevich asks whether periodic tiling billiards have only finitely many generic types of behavior, and in particular whether typical trajectories always fall into a dichotomy between density in open planar sets and integrable behavior, meaning periodicity or linear escape. The excerpt leaves the dynamics, genericity conventions, and several hypotheses undefined, so a precise standalone formulation requires further source clarification.
\subsection{Sources}
\begin{itemize}
\item Local UnsolvedMath registry, permanent ID \(800477\), exact source problem code \texttt{catalogue entry}; supplied title: \emph{Billiards and Quantitative Symplectic Geometry --- Olga Paris-Romaskevich --- Problem 3}. The registry labels the problem \texttt{open}; this is supplied status metadata, not an independently verified mathematical assessment.
\item Supplied source list: Albach, Barraud et al., \emph{Open Problems in Billiards and Quantitative Symplectic Geometry} (2026), source item Olga Paris-Romaskevich, Problem 3. Contributor attribution: Olga Paris-Romaskevich. Supplied URL: \url{https://arxiv.org/abs/2602.12896}. The record reports access on 2026-10-16 and extraction by \texttt{arxiv\_latex}; these are supplied provenance details, not a claim of independent source access or verification.
\item References appearing only as unresolved source labels: [Dynnikov\_main], specifically Theorem 1; [PaRoH]; [BDFI18]; and [TAF]. The excerpt attributes fine Morse theory to Ivan Dynnikov, the initial global-folding ingredient to [BDFI18], and further development to [TAF, PaRoH]. Full bibliographic records and reference-specific links are not supplied.
\end{itemize}
\end{document}
